\documentclass{article}
\usepackage{amsmath,amssymb,hyperref}
\title{Synchronization through an adjoint output comparison}
\date{15 November 2026}
\begin{document}\maketitle
\paragraph{Scope.} A conditional generalization of output-preserving dialgebra synchronization. The mate argument and its orientation have been independently checked. It yields native, possibly noninvertible projection cells, not the pseudo-projection conclusion used in the earlier assembly criterion. No claim that the required adjoint exists universally is made.

\paragraph{Data.} Let $P,Q:L\to K$ be pseudofunctors of bicategories. The native dialgebra has objects $(X,f:PX\to QX)$ and arrows $(h,\alpha:Qh\,f\Rightarrow g\,Ph)$, with the usual compatible base 2-cells. Let a small ordinary diagram of such objects be given. Suppose its underlying diagram $X$ has a chosen lax bilimit $E$ with projections $p_j$ and cone comparisons, and $QX$ has a chosen lax bilimit $Z$ with projections $z_j$. Write $k:QE\to Z$ for the canonical comparison determined by $Qp_j$; its specified cone comparison $z_jk\cong Qp_j$ is invertible. Assume that $k$ has a chosen right adjoint $r:Z\to QE$, with unit and counit satisfying both triangle equations.

The diagram's transitions and native squares form a cone from $PE$ to $QX$, with components $f_jPp_j$. Let $c:PE\to Z$ represent it, with specified invertible cone comparisons $z_jc\cong f_jPp_j$.

\paragraph{Construction.} Give $E$ transition $f_E=rc$. Its projection square is
\[
 Qp_j\,rc\cong z_j krc\xRightarrow{z_j\epsilon c}z_jc\cong f_jPp_j.
\]
Compatibility with the diagram cone cells follows from the cone comparison for $k$, the cone comparison for $c$, and naturality of the counit. The displayed squares need not be invertible.

\paragraph{Universal property.} Fix a source dialgebra $(V,g)$ and a native cone from it to the diagram. The base lax-bilimit equivalence provides $h:V\to E$ and an invertible comparison between its projected base cone and the given base cone. Transport the native cone cells along that comparison. By the $QX$ cone equivalence these cells assemble uniquely into a 2-cell
\[
 \alpha:k\,Qh\,g\Rightarrow c\,Ph.
\]
Here all pseudofunctor and associativity comparisons are retained. Adjunction gives the natural bijection
\[
 K(QV,QE)(Qh\,g,rc\,Ph)
 \cong K(QV,Z)(kQh\,g,c\,Ph).
\]
Its inverse sends $\alpha$ to $r\alpha$ after the unit at $Qh g$. Call the resulting cell $\gamma$. The counit mate equation says exactly that the projected dialgebra arrow $(h,\gamma)$ recovers the given native cone. Thus the induced cone functor is essentially surjective.

For modifications, the underlying base cone equivalence is fully faithful. A candidate base 2-cell between arrows lifts to a dialgebra modification precisely when its two structural pastes agree. Under the mate bijection this is precisely equality of the assembled cone pastes. The $QX$ cone equivalence detects that equality componentwise. Hence the cone functor is fully faithful. All identifications are natural in source arrows and native modifications, using pseudonaturality of the bilimit equivalences and of adjunction hom bijections. Therefore $(E,rc)$ is the desired lax bilimit in the native dialgebra bicategory.

\paragraph{Relationship to preservation.} If $k$ is an equivalence, its adjoint inverse has invertible counit and this recovers the previous pseudo-projection construction. A right adjoint alone is strictly weaker data: the projections may now carry noninvertible cells. Consequently an assembly theorem relying on preservation of pseudo projections cannot be imported unchanged. The new construction establishes existence and the native universal property; compatibility with a particular parallel, state or feedback operation is another theorem.

\paragraph{Application boundary.} Enriched completions supply some underlying and image-diagram cone objects. They do not by themselves supply a right adjoint to every comparison $k$. In particular the old tensor need not become the categorical product. A possible next test is when $k$ arises from a diagonal-index restriction module and the index admits the adjoint needed to give $k$ a right adjoint; this requires exact variance and module verification, not an appeal to an arbitrary companion/conjoint for every original pure map.

\paragraph{Attribution.} The proof uses standard bicategorical mates and representable bilimit universal properties. It is a scoped extension of the project synchronization criterion, not a novelty claim about adjunctions or dialgebras. The completion inputs remain the established KLSS local-cocompletion and Garner--Shulman enriched-module constructions; they are not being cited as proving the additional comparison-adjoint hypothesis.
\end{document}
